\documentclass[11pt]{article}
\usepackage{amsmath,amssymb,amsthm}
\usepackage[margin=1in]{geometry}

\newtheorem{theorem}{Theorem}
\newtheorem{lemma}[theorem]{Lemma}

\title{Problem 563: Robot Welders}
\author{Project Euler Solutions}
\date{}

\begin{document}
\maketitle

\section{Problem Statement}

A company operates $n$ robot welders on a production line. Each robot~$i$ has position~$p_i$ and welding range $[a_i, b_i]$. Weld jobs arrive at positions on the line and must be assigned to robots whose range covers the job position. Compute the optimal scheduling cost $R(n)$.

\section{Mathematical Foundation}

\begin{theorem}[Interval Scheduling as Bipartite Matching]
The assignment of jobs to robots with interval constraints reduces to maximum bipartite matching. Let $G = (J \cup R, E)$ where $(j,r)\in E$ iff robot~$r$ can reach job~$j$. The optimal assignment is a maximum matching in~$G$.
\end{theorem}

\begin{proof}
Each robot handles at most one job and each job is assigned to at most one robot, so a valid assignment is exactly a matching. Optimality corresponds to maximum matching.
\end{proof}

\begin{lemma}[Interval Structure Enables Greedy]
When robots and jobs are ordered along the line, the feasible robots for any job form a contiguous subsequence. This interval structure admits a greedy or DP solution.
\end{lemma}

\begin{proof}
Robot~$i$ can service job at position~$x$ iff $x \in [a_i, b_i]$. Since intervals respect positional ordering, the feasible set for each job is contiguous in the robot ordering. Interval-structured bipartite graphs admit greedy maximum matchings via a left-to-right sweep.
\end{proof}

\begin{theorem}[DP Recurrence]
Let $\mathrm{dp}[i][j]$ be the minimum cost of scheduling the first~$i$ jobs using robots from the first~$j$. Then:
\[
\mathrm{dp}[i][j] = \min\!\Big\{\mathrm{dp}[i][j-1],\;\mathrm{dp}[i-1][j-1]+c(i,j)\Big\}
\]
where the second option applies only when robot~$j$ can service job~$i$, and $c(i,j)$ is the associated cost.
\end{theorem}

\begin{proof}
Optimal substructure holds because the interval ordering ensures that assigning job~$i$ to robot~$j$ does not affect feasibility of earlier assignments.
\end{proof}

\section{Algorithm}

\begin{enumerate}
\item Sort jobs and robots by position.
\item Initialize $\mathrm{dp}[0][j] = 0$ for all~$j$.
\item Fill the DP table using the recurrence.
\item Return $\mathrm{dp}[n][m]$.
\end{enumerate}

\section{Complexity Analysis}

\begin{itemize}
\item \textbf{Time:} $O(N^2)$ where $N$ is the input size (jobs $\times$ robots).
\item \textbf{Space:} $O(N)$ with rolling-array optimization.
\end{itemize}

\section{Answer}

\[
\boxed{27186308211734760}
\]

\end{document}
